\chapter{Родитель связи хопфовских путей с двумя резервуарами}
\label{ch:v10-hopf-path-reservoir-coupling-parent}

\section{Ориентированная задача назначения}

Предыдущий гейт выделил два пути равнорёберного хопфовского цикла,
\begin{equation}
 N_{\rm path}=\operatorname{diag}(1,2),\qquad F_e=\log2,
\end{equation}
но не назначил их горячей и холодной ваннам. Центрируем оператор длин и
введём уже типизированную ориентацию резервуаров:
\begin{equation}
 Z_{\rm path}=2N_{\rm path}-3I=\operatorname{diag}(-1,1),\qquad
 Z_{\rm res}=\operatorname{diag}(-1,1).
\end{equation}
Для перестановочного морфизма $U$ измерим дефект
\begin{equation}
 D(U)=Z_{\rm res}-UZ_{\rm path}U^*.
\end{equation}

\section{Знак разности устраняет перестановочную неоднозначность}

Для согласованного назначения $U=I$ имеем $D(I)=0$. Единственная другая
перестановка $U=X$ меняет ориентацию:
\begin{equation}
 XZ_{\rm path}X=-Z_{\rm path},\qquad
 D(X)=\operatorname{diag}(-2,2).
\end{equation}
Соответствующие квадраты дефекта равны $0$ и $8$. Более того,
\begin{align}
 U=I &: (a_h,a_c)=(\log2,2\log2),&a_c-a_h&=\log2,\\
 U=X &: (a_h,a_c)=(2\log2,\log2),&a_c-a_h&=-\log2.
\end{align}
Поэтому ранее выведенный положительный знак разности единственно выбирает
назначение «горячая ванна --- путь длины один, холодная --- путь длины два».

\section{Условный родитель и отсутствующий смешанный блок}

Если морфизм назначения допущен, минимальный родитель
\begin{equation}
 S_{\rm assign}=\frac12\left[(a_h-\log2)^2+(a_c-2\log2)^2\right]
\end{equation}
имеет гессиан $I_2$, ранг $2$ и определитель $1$. Он строго выбирает оба
сродства и даёт KMS-отношения $(1/2,1/4)$.

Однако унаследованный родитель остаётся прямой суммой резервуарного и
хопфовского секторов. Его смешанный блок, способный породить $U$, равен
нулю. Поэтому найдено уникальное условное назначение, но не происхождение
самого межсекторного морфизма.

Наконец, равенства $a_r=\beta_r\Delta$ сохраняют ядро
\begin{equation}
 (\log\beta_h,\log\beta_c,\log\Delta)
 \longmapsto
 (\log\beta_h-s,\log\beta_c-s,\log\Delta+s).
\end{equation}
Температурная карта имеет ранг $2$ и одномерное ядро $(-1,-1,1)$.

\begin{theorem}[Уникальность ориентированного путевого назначения]
При фиксированных длинах $1:2$, силе ребра $\log2$ и выведенном знаке
$a_c-a_h=\log2$ единственным перестановочным морфизмом является
$U=I$: горячей ванне назначается путь длины один, холодной --- длины два.
\end{theorem}

\begin{nogocriterion}[Уникальная совместимость не порождает связь]
Нулевой дефект условного морфизма не считается его физическим происхождением,
пока унаследованный общий родитель имеет нулевой резервуарно-хопфовский
смешанный блок. Без независимой энергетической щели он также не задаёт
абсолютные температуры.
\end{nogocriterion}

\begin{protocol}\begin{enumerate}
\item условная архитектура: \textbf{$4/4$};
\item ориентированное назначение при заданной связи: \textbf{$1/1$};
\item происхождение смешанного морфизма: \textbf{$0/1$};
\item физическая температура: \textbf{$0/1$};
\item итоговый origin-ledger: \textbf{$4/6$};
\item следующий гейт: \textbf{аудит кандидатов резервуарно-хопфовского интертвинера}.
\end{enumerate}\end{protocol}

Машинный сертификат сохранён в
\path{s2t_v10_cell_birth_four_volume_induced_newton_breathing_anomaly_hopf_path_length_reservoir_coupling_parent_origin_gate_results.json}.